\uuid{AHfa}
\titre{Application linéaire}
\auteur{hanani ; mzari}
\organisation{exo7}
\licence{CC BY-NC-SA 4.0}
\datecreate{2020-01-20}
\niveau{L1}
\module{Algèbre}
\chapitre{Application linéaire}
\sousChapitre{Définition}
\difficulte{1}
\theme{Applications linéaires, Application linéaire}
\exo7qcmid{lille-2.applications-lineaires.003}

\contenu{
\begin{question}
\qclassification{107.01}
\qtags{motcle=Application linéaire}
On considère les deux applications suivantes :
$$\begin{array}{rccc}f:&\Rr^3&\to&\Rr^2\\
& (x,y,z)&\to &(x+y,x-z) \end{array} \quad \mbox{et} \quad \begin{array}{rccc}g:&\Rr^3&\to&\Rr^2\\
& (x,y,z)&\to &(xy,xz). \end{array}$$
Quelles sont les assertions vraies ?
\begin{answers}
\good{$f(0,0,0)=(0,0)$}.
\good{$f$ est une application linéaire}.
\bad{$g(1,1,0)=g(1,0,0)+ g(0,1,0)$}.
\bad{$g$ est une application linéaire}.
\end{answers}
\begin{explanations} $f$ est linéaire. $g$ ne l'est pas,
puisque $g((1,0,0)+ (0,1,0)) = g(1,1,0) = (1,0)$ et $ g(1,0,0)+ g(0,1,0) = (0,0)$.
\end{explanations}
\end{question}
}
